\chapter{The Intermediate Representation}\label{r01:ch:ir}

All backend passes operate on one intermediate representation (IR). It consists of
\emph{objects}, which are symbols, buffers, constants, literals and functions, connected
by an \emph{abstract syntax tree} (AST) of statements and expressions. Every object and
every expression node carries a data type from the \ps{} type system. This chapter
introduces these building blocks. The upstream developer documentation
(\file{docs/source/backend/}) describes them in more depth~\cite{pystencils-docs}.

\section{The kernel creation context}

A \cls{KernelCreationContext} (\file{backend/kernelcreation/context.py}) holds all state
that is shared by the passes of one kernel:
\begin{itemize}
  \item the \textbf{default types}: \code{default_dtype} and \code{index_dtype};
  \item the \textbf{symbol table}, which maps names to \cls{PsSymbol} objects.
        \code{get_symbol(name, dtype)} finds or creates a symbol,
        \code{get_new_symbol(name)} creates one with a fresh name,
        \code{duplicate_symbol(s)} creates a copy named \code{s__0}, \code{s__1}, \dots;
  \item the \textbf{fields} of the kernel, grouped into domain, index, buffer and custom
        fields, and for each field its \cls{PsBuffer} (\code{get_buffer(field)});
  \item the \textbf{iteration space} (\cref{r01:sec:ispace});
  \item \textbf{reduction info}: for each reduction symbol a write-back pointer, a local
        accumulator and the initial value;
  \item \textbf{required headers} and free-form \textbf{metadata}, which ends up in
        \code{Kernel.metadata}.
\end{itemize}

\pysnippet{ex16_ir.context}
\outsnippet{ex16_ir.context}

\section{IR objects}

\paragraph{Memory model.}
To reason about aliasing, the backend distinguishes three kinds of memory, which never
overlap:
\begin{itemize}
  \item \textbf{symbols} act as virtual registers: private, never aliased, never
        addressed;
  \item \textbf{field buffers} are disjoint arrays (the generated code declares their
        pointers \code{RESTRICT});
  \item the \textbf{unmanaged heap} is where all other raw pointers point. Any two raw
        pointers are assumed to alias.
\end{itemize}

\paragraph{Symbols.}
\cls{PsSymbol} (\file{backend/memory.py}) has a name, a data type (possibly \code{None}
until typification) and a set of \term{properties}. Symbols are mutable and compared by
identity. Properties attach meaning to a symbol. \code{FieldShape}, \code{FieldStride}
and \code{BufferBasePtr} mark the size, stride and base-pointer symbols of buffers.
\code{BackendPrivateProperty} marks properties that are not exported to the kernel's
parameters. When the kernel is finalized, the properties of parameter symbols become the
properties of the corresponding \code{Parameter} (\cref{r01:sec:kernel-object}).

\paragraph{Buffers.}
\cls{PsBuffer} describes the memory of a field: name, element type, a base pointer
symbol, and a list of shape and stride entries. Each entry is either a symbol or a
constant. Index dimensions are always present in the buffer: a scalar field gets a
trailing index dimension of size 1, so a 2D scalar field has a 3D buffer. The buffer
above belongs to \code{f(3): float32[2D]}. Its last shape entry is the constant 3, and
the base pointer \code{_data_f} carries the \code{BufferBasePtr} property.

\paragraph{Constants and literals.}
\cls{PsConstant} (\file{backend/constants.py}) wraps a number together with an optional
type. An untyped constant receives its type during typification
(\code{interpret_as(dtype)}). \cls{PsLiteral} (\file{backend/literals.py}) is a piece of
C code with a type, printed verbatim. Examples are \code{blockIdx.x} or
\code{PS_FP64_INFINITY}.

\paragraph{Functions.}
\file{backend/functions.py} distinguishes two kinds of functions:
\begin{itemize}
  \item \textbf{IR functions} (\cls{PsIrFunction}) are abstract operations that a
        platform must lower: \cls{PsMathFunction} (\code{exp}, \code{sqrt},
        \code{min}, \dots, listed in \code{MathFunctions}), \cls{PsConstantFunction}
        ($\pi$, $e$, $\pm\infty$), \cls{PsGpuIndexingFunction}
        (\code{threadIdx.x}, \dots), \cls{PsGpuIntrinsicFunction}, \cls{PsRngEngineFunction}
        (Philox) and \cls{PsReductionWriteBack}.
  \item \textbf{C functions} (\cls{CFunction}) have a concrete name and signature. They
        are the only functions allowed in a kernel when it is printed. Vector
        intrinsics such as \code{_mm256_add_pd} or \code{__riscv_vfadd_vv_f64m1} are
        \code{CFunction}s created by the platforms.
\end{itemize}

\section{The type system}\label{r01:sec:types}

Types are instances of \cls{PsType} (\file{src/pystencils/types/}). The class hierarchy
is shown in \cref{r01:fig:types}. A metaclass interns all type objects, so
\code{create_type("float64") is Fp(64)} holds, and types can be compared with \code{is}
or \code{==}. Every type has a \code{const} flag. Pointer types also have a
\code{restrict} flag. \code{constify(t)} and \code{deconstify(t)} switch between the two
forms. \code{c_string()} gives the C spelling.

\begin{figure}[htbp]
  \centering
  % The PsType class hierarchy (src/pystencils/types/types.py)
\begin{forest}
  for tree={
    font=\ttfamily\footnotesize, draw=psblue!60, fill=psblue!5, rounded corners=1.5pt,
    inner sep=2.5pt, grow'=east, parent anchor=east, child anchor=west,
    anchor=west, l sep=6mm, s sep=1.2mm, edge={draw=psgray},
    edge path={\noexpand\path[\forestoption{edge}] (!u.parent anchor) -- +(2.5mm,0)
      |- (.child anchor)\forestoption{edge label};},
  }
  [PsType
    [PsCustomType]
    [PsVoidType]
    [PsStructType]
    [PsDereferencableType
      [PsPointerType]
      [PsArrayType
        [PsNamedArrayType [PsShortArrayType]]]]
    [PsNumericType
      [PsVectorType]
      [PsScalarType
        [PsBoolType]
        [PsIntegerType [PsSignedIntegerType] [PsUnsignedIntegerType]]
        [PsIeeeFloatType]]]
  ]
\end{forest}

  \caption[The \texttt{PsType} hierarchy]{The \texttt{PsType} hierarchy.}
  \label{r01:fig:types}
\end{figure}

\code{create_type} parses user specifications: type names, NumPy types, Python types
and \code{PsType} objects. The module \code{pystencils.types.quick} provides short aliases
(\code{Fp(64)}, \code{SInt(32)}, \code{UInt(8)}, \code{Ptr(...)}, \code{Arr(...)},
\code{Bool()}, \code{Custom("std::vector<int>")}), which platforms also use in
\code{match} statements.
\pysnippet{ex16_ir.types}
\outsnippet[nolol=false,caption={User type specifications, their \texttt{PsType}, IR spelling and C spelling.},basicstyle=\ttfamily\scriptsize]{ex16_ir.types-table}
Each scalar type knows how to print a typed literal. This is why the generated code
contains \code{1LL} and \code{1.0f}:
\outsnippet{ex16_ir.literals}
\cls{PsVectorType} (e.g.\ \code{float64<4>}) exists only inside the IR. It has no C
spelling, because each vector platform maps it to its own intrinsic type
(\code{__m256d}, \code{float64x2_t}, \code{vfloat64m1_t}, \dots). \cls{PsStructType}
represents NumPy structured dtypes, such as the elements of index fields.
\cls{PsCustomType} wraps arbitrary C++ type names.

\section{The abstract syntax tree}\label{r01:sec:ast}

AST nodes (\file{backend/ast/}) derive from \cls{PsAstNode}, which provides
\code{children}, \code{clone()} and \code{structurally_equal()}. \code{str(node)} prints
the node as IR. There are four groups of nodes:

\begin{description}
  \item[Structural nodes] (\file{structural.py}): \cls{PsBlock}, \cls{PsStatement},
    \cls{PsAssignment}, \cls{PsDeclaration} (an assignment that declares its symbol),
    \cls{PsLoop}, \cls{PsConditional}, \cls{PsPragma}, \cls{PsComment}.
  \item[Expressions] (\file{expressions.py}): leaves \cls{PsSymbolExpr},
    \cls{PsConstantExpr} and \cls{PsLiteralExpr}; memory accesses \cls{PsBufferAcc}
    (multidimensional, on a buffer), \cls{PsMemAcc} (pointer plus offset),
    \cls{PsSubscript} (arrays) and \cls{PsLookup} (struct members); \cls{PsCall},
    \cls{PsTernary}, \cls{PsCast}, \cls{PsAddressOf}; arithmetic, integer, boolean and
    relational operators (\cls{PsAdd}, \cls{PsIntDiv}, \cls{PsBitwiseAnd}, \cls{PsAnd},
    \cls{PsLt}, \dots); \cls{PsArrayInitList}.
  \item[Vector nodes] (\file{vector.py}): \cls{PsVecBroadcast} (scalar to all lanes),
    \cls{PsVecMemAcc} (contiguous or strided vector load/store) and
    \cls{PsVecHorizontal} (horizontal reduction of a vector into a scalar).
  \item[Iteration axes] (\file{axes.py}): \cls{PsAxesCube}, \cls{PsAxisRange} and the
    \cls{PsIterationAxis} subclasses \cls{PsLoopAxis}, \cls{PsParallelLoopAxis},
    \cls{PsSimdAxis} and the GPU axes. They describe how to iterate before any loop
    exists (\cref{r01:sec:axes}).
\end{description}

The following snippet builds a small AST by hand with the \cls{AstFactory}. The SymPy
assignments go through freeze and typify (\cref{r01:ch:translation}), and \code{loop} wraps
them in a loop:
\pysnippet{ex16_ir.ast}
\irsnippet[nolol=false,caption={The hand-built AST printed as IR (\texttt{emit\_ir}).}]{ex16_ir.ast-ir}
\outsnippet[nolol=false,caption={The same AST printed as C (\texttt{emit\_code}).},style=cpp]{ex16_ir.ast-c}
The tree structure of the first declaration, with the type of every node:
\outsnippet{ex16_ir.tree}

\paragraph{Canonical form.}
Most transformations require their input to be in \term{canonical form}:
\begin{enumerate}
  \item every symbol, constant and expression has a data type;
  \item every symbol is declared at most once;
  \item a symbol that is never written after its declaration has a \code{const} type;
  \item a non-\code{const} symbol has at least one assignment besides its declaration.
\end{enumerate}
The \cls{Typifier} establishes rule~1, and \cls{CanonicalizeSymbols} the other three.
The hand-built AST above is not yet canonical: \code{a} and \code{c} are not
\code{const}, although they are never reassigned. In real kernels the driver runs
\cls{CanonicalizeSymbols}, which is why the declarations in \cref{r01:lst:heat-code} are
\code{const}.

\paragraph{Traversal and analysis helpers.}
\file{backend/ast/iteration.py} provides \code{dfs_preorder} and \code{dfs_postorder}.
\file{backend/ast/analysis.py} provides \code{collect_undefined_symbols} (the future
parameters), \code{collect_required_headers} and an \cls{OperationCounter}.
\file{backend/transformations/rewrite.py} contains \code{substitute_symbols} and
\code{collapse_blocks}.

\section{Printers}

\file{backend/emission/} contains two printers that share a base class with C operator
precedence rules:
\begin{itemize}
  \item \cls{IRAstPrinter} (\code{emit_ir}) prints every node, including vector and axes
        nodes, with type annotations (\code{[1.0: float64]}, \code{ctr_0: int64}).
  \item \cls{CAstPrinter} (\code{emit_code}) prints valid C. It refuses nodes that must
        have been lowered first (\cls{PsBufferAcc}, \cls{PsVecMemAcc}), and prints
        literals through the types' \code{create_literal}.
\end{itemize}
For a \code{Kernel}, \code{print_signature} prints \code{void <name>(<params>)} with a
\code{__global__} prefix for GPU kernels.
